% EN48 c00 | 原书 PDF 1–9（原书 i–v） | 前置页：摘要、献词、致谢（封面/版权页信息并入扉页，目录与图目录由 LaTeX 生成）
% 骨架文件：标题行已设定，正文由译者填写。每页起首用 % ===== FRONT MATTER (PDF n) ===== 注释标记。
% ===== FRONT MATTER (PDF 4) =====
\section*{摘要}
\addcontentsline{toc}{section}{摘要}

\begin{center}
{\small
亏格 2 曲线的同调镜像对称及其广义 Strominger--Yau--Zaslow 镜像\\[0.4em]
Catherine Kendall Asaro Cannizzo\\
数学哲学博士，University of California, Berkeley\\
导师：Denis Auroux 教授（主席）
}
\end{center}

受物理学中的观察启发，镜像对称是说某些流形成对出现：$X$ 与 $Y$，且 $X$ 上的复几何与
$Y$ 上的辛几何互为镜像，于是可以由 $X$ 的已知性质推断 $Y$ 的信息。
Strominger--Yau--Zaslow（1996）指出，这类配对在几何上来自具有同一底空间、纤维相互关联的
环面丛，称为 SYZ 镜像对称。Kontsevich（1994）猜想：$X$ 上的一个复不变量（有界相干层
导出范畴）应当等价于 $Y$ 上的一个辛不变量（深谷范畴）。此即同调镜像对称。本项目先用
Abouzaid--Auroux--Katzarkov 关于阿贝尔簇中超曲面的 SYZ 镜像构造，得到作为流形的 $X$
与 $Y$。复流形一侧来自亏格 2 曲线作为其 Jacobi 环面中的超曲面；我们给 SYZ 镜像流形
配上一个辛形式。随后描述复一侧的范畴到 $Y$（视为辛纤维丛）的一个上同调
Fukaya--Seidel 范畴的嵌入。本文的纤维丛属于最早配备深谷范畴的非正合、非 Lefschetz
纤维丛之列，但定义它的主要几何想法与 Seidel 对 Lefschetz 纤维丛的深谷范畴的构造相同。

\transnote{摘要里的 ``bounded derived category of coherent sheaves'' 即正文的
$\DBCoh$，本译本统一作「有界相干层导出范畴」。注意此处的 ``embedded'' 是范畴嵌入，
不是流形嵌入：$H$ 本身是 $V$ 中的超曲面，这里说的是 $\DBCoh(H)$ 嵌入 $Y$ 的
Fukaya--Seidel 范畴。}

\transnote{``cohomological Fukaya--Seidel category'' 指把态射复形取上同调后得到的
普通范畴，记作 $H^0(\mathrm{FS}(Y,v_0))$；它丢掉了 $A_\infty$ 的高阶运算，
所以第 4、5 章证明的只是上同调层面的等价，而非 $A_\infty$-等价。}

\transnote{``nonexact'' 指辛形式 $\omega$ 不正合：$Y$ 是带有紧纤维的辛纤维丛，
纤维体积非零，故 $\omega$ 不是恰当形式。这与 Seidel 在 [Sei08] 中处理的
Lefschetz 纤维丛（正合情形）有实质差别。}

\transnote{``Jacobian torus'' 即 $\mathrm{Pic}^0(\Sigma_2)$，见 \S 1.2：
固定点 $x_0$ 后，$x\mapsto\OO(x-x_0)$ 把 $\Sigma_2$ 嵌入这个二维阿贝尔簇。
「Jacobi 环面」与「阿贝尔簇」在本文指同一个对象的不同侧面，前者强调它是复环面。}

% ===== FRONT MATTER (PDF 5) =====
\section*{献词}
\addcontentsline{toc}{section}{献词}

\begin{center}
谨以本论文献给我的父亲 John Kendall Cannizzo，献上我全部的爱。

感谢您做我生命中的阳光。愿您安息。
\end{center}

\transnote{献词中的人名保留英文原名，不音译。这是全书的通例：凡作者、导师、
被致谢者，一律照录原文。}

% ===== FRONT MATTER (PDF 9) =====
\section*{致谢}
\addcontentsline{toc}{section}{致谢}

本工作部分受 NSF 基金 DMS-1264662、DMS-1406274、DMS-1702049 与 Simons Foundation
基金（\# 385573，Simons 同调镜像对称合作项目）资助。感谢我的导师 Denis Auroux
在本论文项目上无法估量的学术支持与指导。感谢 Mohammed Abouzaid、Melissa Liu、
Katrin Wehrheim、Sheel Ganatra、Heather Lee、Roberta Guadagni、Jingyu Zhao、
Sara Venkatesh、Wolfgang Schmaltz、Zhengyi Zhou、Benjamin Filippenko，以及 MSRI
导出范畴项目 [APS] 与我的有益数学讨论。感谢 Jo Nelson、Ailsa Keating、Bahar Acu、
Ziva Myer、Yu Pan、Morgan Weiler 与 Anastasia Chavez 在申请过程与职业发展上给予的
宝贵建议与支持。感谢 Noetherian Ring 与 UC Berkeley 数学系的工作人员。最后，
感谢我的父母始终不渝的支持。

\transnote{``Noetherian Ring'' 是 UC Berkeley 数学系的研究生学术组织，不是环论
概念。这里按专名保留英文，避免误读成「诺特环」。}

\transnote{MSRI 即 Mathematical Sciences Research Institute（伯克利）。
``Derived Categories program'' 指 2018 年前后 MSRI 的导出范畴专题项目，
方括号 [APS] 是原文的引用标记，指参与该项目的 Abramovich--Polishchuk--Segal
等人在相关主题上的工作；详见参考文献。}

\transnote{致谢中的基金编号照录原文。``Simons Collaboration on Homological
Mirror Symmetry'' 是 Simons Foundation 资助的协作项目，不是期刊或会议。}
